\documentclass[12pt]{article}
\usepackage{amsmath, amssymb}
\usepackage{geometry}
\geometry{margin=1in}
\usepackage{enumitem}

\title{CSE400 – Fundamentals of Probability in Computing}
\author{Lecture 8: Gaussian, Uniform, Exponential, and Gamma Random Variables}
\date{}

\begin{document}
\maketitle

\section*{Outline}

\begin{itemize}
\item Gaussian Random Variable
\begin{itemize}
\item Definition and Properties
\item Standard Forms: Q-function and Error Function
\item Relation between $\Phi$-function and Q-function
\item Example
\end{itemize}
\item Types of Continuous Random Variable
\item Uniform Random Variable
\item Exponential Random Variable
\item Laplace Random Variable
\item Gamma Random Variable
\item Graph and Special Cases
\item Example
\item Homework Problem
\item Problem Solving
\item In-Class Activity: Gaussian Density Estimation
\end{itemize}

\section*{Moments and Central Moments}

For a continuous random variable $X$:

\[
E[(X-\mu_X)^n] = \int (x-\mu_X)^n f_X(x)\, dx
\]

For a discrete random variable:

\[
E[(X-\mu_X)^n] = \sum_k (x_k-\mu_X)^n P_X(x_k)
\]

\subsection*{Case $n=0$}

\[
E[(X-\mu_X)^0] = E[1] = 1
\]

\subsection*{Case $n=1$}

\[
E[(X-\mu_X)] = E[X] - \mu_X = \mu_X - \mu_X = 0
\]

\subsection*{Case $n=2$ (Variance)}

\[
\sigma_X^2 = E[(X-\mu_X)^2]
\]

Expand:

\[
= E[X^2 - 2\mu_X X + \mu_X^2]
\]

\[
= E[X^2] - 2\mu_X E[X] + \mu_X^2
\]

\[
= E[X^2] - \mu_X^2
\]

\subsection*{Case $n=3$: Skewness}

\[
C_s = \frac{E[(X-\mu_X)^3]}{\sigma_X^3}
\]

\[
C_s > 0 : \text{PDF is right skewed}
\]

\[
C_s < 0 : \text{PDF is left skewed}
\]

\subsection*{Case $n=4$: Kurtosis}

\[
C_k = E[(X-\mu_X)^4]
\]

Large value implies a large peak near the mean.

\section*{Linearity of Expectation}

For constants $a$ and $b$:

\[
E[aX+b] = aE[X] + b
\]

For a sum of functions:

\[
E\left[\sum_{k=1}^N g_k(X)\right] = \sum_{k=1}^N E[g_k(X)]
\]

Also:

\[
E[X] = \int x f_X(x)\, dx
\]

\section*{Gaussian Random Variable}

\subsection*{Definition}

A Gaussian random variable is one whose PDF can be written as:

\[
f_X(x) = \frac{1}{\sqrt{2\pi\sigma^2}} \exp\left(-\frac{(x-m)^2}{2\sigma^2}\right)
\]

where

\[
m = \mu_X
\]

\[
\sigma^2 \text{ is the variance}
\]

Notation:

\[
X \sim \mathcal{N}(m,\sigma^2)
\]

\section*{Standard Forms}

\subsection*{Error Function}

\[
\text{erf}(x) = \frac{2}{\sqrt{\pi}} \int_0^x e^{-t^2} dt
\]

\subsection*{Complementary Error Function}

\[
\text{erfc}(x) = 1 - \text{erf}(x)
\]

\[
= \frac{2}{\sqrt{\pi}} \int_x^{\infty} e^{-t^2} dt
\]

\subsection*{$\Phi$-function (CDF of Standard Normal)}

\[
\Phi(x) = \frac{1}{\sqrt{2\pi}} \int_{-\infty}^x \exp\left(-\frac{t^2}{2}\right) dt
\]

\subsection*{Q-function (Gaussian Tail Function)}

\[
Q(x) = \frac{1}{\sqrt{2\pi}} \int_x^{\infty} \exp\left(-\frac{t^2}{2}\right) dt
\]

\subsection*{Identity}

\[
Q(x) = 1 - \Phi(x)
\]

\section*{Relation Between $\Phi$-function and Q-function}

\subsection*{Evaluating the CDF}

\[
F_X(x) = \int_{-\infty}^x \frac{1}{\sqrt{2\pi\sigma^2}} 
\exp\left(-\frac{(y-m)^2}{2\sigma^2}\right) dy
\]

Substitute:

\[
t = \frac{y-m}{\sigma}
\]

Then:

\[
F_X(x) = \int_{-\infty}^{\frac{x-m}{\sigma}} \frac{1}{\sqrt{2\pi}}
\exp\left(-\frac{t^2}{2}\right) dt
\]

\[
F_X(x) = \Phi\left(\frac{x-m}{\sigma}\right)
\]

\subsection*{Evaluating Tail Probability}

\[
P(X>x) = \int_{\frac{x-m}{\sigma}}^{\infty} \frac{1}{\sqrt{2\pi}}
\exp\left(-\frac{t^2}{2}\right) dt
\]

\[
P(X>x) = Q\left(\frac{x-m}{\sigma}\right)
\]

\subsection*{Key Identities}

\[
Q(x) = 1 - \Phi(x)
\]

\[
F_X(x) = 1 - Q\left(\frac{x-m}{\sigma}\right)
\]

\section*{Notes from Slides}

To evaluate the CDF of a Gaussian RV, evaluate the $\Phi$-function at the point.

The Q-function is more natural for evaluating probabilities of the form $P(X>x)$.

\section*{Example}

Given:

\[
f_X(x) = \frac{1}{\sqrt{8\pi}} \exp\left(-\frac{(x+3)^2}{8}\right)
\]

\subsection*{Comparison with Gaussian PDF}

Standard Gaussian PDF:

\[
f_X(x) = \frac{1}{\sqrt{2\pi\sigma^2}} \exp\left(-\frac{(x-m)^2}{2\sigma^2}\right)
\]

Identifying parameters:

\[
m = -3
\]

\[
\sigma^2 = 4 \Rightarrow \sigma = 2
\]

\subsection*{Gaussian Relations Used}

\[
F_X(x) = \Phi\left(\frac{x-m}{\sigma}\right)
\]

\[
P(X>x) = Q\left(\frac{x-m}{\sigma}\right)
\]

\subsection*{Required Probabilities}

\[
P(X \le 0)
\]

\[
P(X > 4)
\]

\[
P(|X+3| < 2)
\]

\[
P(|X-2| > 1)
\]

\subsection*{Hint (as given)}

Use $Q(x)$ for right-tail probabilities.

Use $\Phi(x)$ for left-tail probabilities.

\end{document}

